\chapter{Results}
\label{ch:results}

\section{Edge balancing}

The edge balancing works reliably. With the automatic trim, the wheel oscillates around zero speed
and the trim settles within about \SI{\pm 0.1}{\degree}. The cube also recovers from gentle pushes,
as shown in the animation in the repository (\cref{fig:edge_results}).

\begin{figure}[htbp]
  \centering
  \begin{subfigure}[b]{0.38\textwidth}
    \centering
    \includegraphics[width=\textwidth]{edge_balance_front_view.jpg}
    \caption{Front view of the cube balancing on its edge.}
    \label{fig:edge_front}
  \end{subfigure}\hspace{0.06\textwidth}
  \begin{subfigure}[b]{0.38\textwidth}
    \centering
    \includegraphics[width=\textwidth]{edge_balance_push_resistance_frame.png}
    \caption{Frame of the animation in which the cube is pushed while balancing on its edge.}
    \label{fig:edge_push}
  \end{subfigure}
  \caption{Balancing on the edge.}
  \label{fig:edge_results}
\end{figure}

\section{Vertex balancing}

After the improvements described in \cref{sec:history}, and in particular after calibrating the
vertex by hand, the cube balances on its vertex for several minutes in a row. Only a small residual
oscillation remains, which is attributed to the limited motor power and to a calibration that is
not perfect. \Cref{fig:vertex_results} shows frames of the two animations in the repository.

\begin{figure}[htbp]
  \centering
  \begin{subfigure}[b]{0.38\textwidth}
    \centering
    \includegraphics[width=\textwidth]{vertical_balance_loop_frame.png}
    \caption{Frame of the balancing loop animation.}
    \label{fig:vertex_loop}
  \end{subfigure}\hspace{0.06\textwidth}
  \begin{subfigure}[b]{0.38\textwidth}
    \centering
    \includegraphics[width=\textwidth]{vertical_balance_around_view_frame.png}
    \caption{Frame of the animation with a view around the cube.}
    \label{fig:vertex_around}
  \end{subfigure}
  \caption{Balancing on the vertex.}
  \label{fig:vertex_results}
\end{figure}

\section{Motor test}

The motor test (\cmd{motortest}, PWM 150 for \SI{200}{ms}) gave encoder counts of 379, 377 and 380
for the motors X, Y and Z, confirming that the three motors deliver the same power, and the
measured gyroscope reactions confirmed the motor signs $(+1, -1, -1)$ and the encoder signs
$(+1, +1, +1)$.

\section{Development history and lessons learned}
\label{sec:history}

The vertex balancing did not work at the first attempt. \Cref{tab:history} summarises the most
important problems, their causes and the resulting balancing times on the vertex. Several of these
lessons apply to any balancing robot.

\begin{longtable}{>{\raggedright\arraybackslash}p{0.27\textwidth}>{\raggedright\arraybackslash}p{0.45\textwidth}>{\raggedright\arraybackslash}p{0.18\textwidth}}
  \caption{Development history of the vertex balancing.}
  \label{tab:history}\\
  \toprule
  Problem & Cause and solution & Vertex time \\
  \midrule
  \endfirsthead
  \toprule
  Problem & Cause and solution & Vertex time \\
  \midrule
  \endhead
  Vertex filter drifted to the wrong sign after about \SI{0.3}{s} &
    The tilt error was computed as $\vect{r}\times\vect{n}$ instead of $\vect{n}\times\vect{r}$,
    so gyroscope and accelerometer had opposite signs. Additionally, the filter used a wrong time
    step and included the yaw rotation. All three were fixed. & -- \\
  Cube slowly tipped away &
    The encoder signs of motors Y and Z were wrong, which cancelled half of the restoring
    action. Fixed with the motor test. & 1--2\,s \\
  Slow rocking with a period of about \SI{1.3}{s} &
    Gains changed from 70\,/\,8\,/\,0.35 to 90\,/\,11\,/\,0.30. & -- \\
  Motor Z stopped working in two runs &
    Loose contact (encoder showed zero despite a large command). & -- \\
  Run stopped by the shock detection although the cube was upright &
    Fast reversals of the wheels produced short accelerations above \SI{2}{g}; stopping and braking
    then threw the cube over. Now only a shock lasting longer than \SI{150}{ms} stops the run;
    \cmd{vk2} reduced to 9.5. & -- \\
  Cracking noises &
    A wheel rubbed against the frame because of a loose screw. & 7--8\,s \\
  Runs ended with the wheels saturated, cube visibly rotating about its vertical axis &
    The wheel momentum about the vertical axis was fed back like the tilt momentum (positive
    feedback). The yaw axis is now controlled separately (\cref{sec:yaw_physics}). & -- \\
  All run times suddenly much shorter (about \SI{2}{s}) &
    Bluetooth output blocked the control loop. Every log line is now sent as one packet, after the
    motor commands. Afterwards, the measured control gap was a constant \SI{6}{ms}. & 7--10\,s \\
  Results varied strongly between sessions &
    The vertex pose had been calibrated in a fixture that held the cube geometrically straight,
    but the centre of mass was slightly off the geometric axis. Calibrating by hand, at the real
    balance point, solved the problem. & several minutes \\
  \bottomrule
\end{longtable}

The main lessons are:
\begin{itemize}
  \item \textbf{Signs and conventions first.} Most early failures were sign errors. A motor test
        that measures the reaction of the cube to each motor saves a lot of time.
  \item \textbf{The balance point is defined by the centre of mass, not by the geometry.}
  \item \textbf{Look for positive feedback} when a quantity grows exponentially.
  \item \textbf{Logging must never block the control loop.}
  \item \textbf{Change only one thing at a time.} In one test series, two changes were made at
        once, and their effects could not be separated.
  \item \textbf{Mechanical problems look like control problems.} A rubbing wheel or a loose
        contact can ruin otherwise good gains.
\end{itemize}
